\section{An exact test for the local first-layer hypothesis}
\label{sec:cm-exact-first-layer}
Let \(p\ge5\), \(h=(p-1)/2\), and let \(g\) be a primitive residue
modulo \(p\).  Put
\[
 \begin{gathered}
 s=[g],\quad c=s^h,\quad O=\mathbb Z[\zeta_p]\otimes\mathbb Z_2,\\
 T=\{u\in1+2O:u\,c(u)=1\},\quad P=(s-1)T,\quad
 \mathscr M=(O/2O)^c.
 \end{gathered}
\]
Operators on unit groups are written additively, so
\((s-1)u=s(u)u^{-1}\).
Let \(Q\) be the integral CM operator, assume \(Q(1)\) is odd, and
assume
\[
 D=\det\bigl(Q\mid\mathbb Z_2[s]/(s^h+1)\bigr)\ne0.
\]
Write
\[
 d=\operatorname{ord}_p(2),\quad r=(p-1)/d,\quad
 e=\begin{cases}r,&d\text{ even},\\ r/2,&d\text{ odd},\end{cases}
\]
and
\[
 b=\deg\gcd(\overline Q,X^h-1),\qquad
 t=\deg\gcd(\overline Q,X^e-1).
\]

\begin{proposition}\label{prop:cm-exact-first-layer}
With every finite sign subgroup included,
\[
 [P:QP]=[T:QT]=2^{v_2(D)+t},\qquad v_2(D)\ge b.
\]
The natural reduction
\[
 \alpha:T\longrightarrow \mathscr M,\qquad
 \alpha(u)=\frac{u-1}{2}\pmod{2O}
\]
induces a surjection \(T/QT\to \mathscr M/\overline QM\).  It is an isomorphism
if and only if \(t=0\) and \(v_2(D)=b\).
The logarithmic convention
\[
 \alpha_{\log}(u)=\frac{\log u}{2}\pmod{2O}
\]
induces a map between the same quotients and has exactly the same
isomorphism criterion.
\end{proposition}

\begin{proof}
The algebra \(O\) is a product of \(r\) unramified degree-\(d\) integer
rings, because the irreducible factors of \(\Phi_p\) modulo \(2\)
correspond to the orbits of multiplication by \(2\) on
\((\mathbb Z/p)^\times\).
If \(d\) is even, \(-1\in\langle2\rangle\), so \(c\) fixes every factor
and acts nontrivially within it.  If \(d\) is odd, \(c\) exchanges the
factors in pairs.  In either case there are exactly \(e\) conjugation
orbits, and \(s\) permutes those orbits cyclically.

The logarithm converges on \(1+2O\), and it is inverse to the exponential
on \(1+4O\) and \(4O\).
If \(\log u=0\), then \(u^2\in1+4O\) has zero logarithm and hence
\(u^2=1\).  Conversely every component sign has logarithm zero.
The norm-one condition says that the signs agree on each conjugation
orbit.  Thus
\[
 \mathcal T:=\operatorname{Tor}(T)=(C_2)^e,
 \qquad \mathcal T\simeq\mathbb F_2[X]/(X^e-1)
\]
as an \(s\)-module.  In particular
\[
 |\mathcal T/Q\mathcal T|=2^t.
\]

The image \(\Lambda=\log T\) is a full lattice in
\(O^-=\{a\in O:c(a)=-a\}\): it is contained in \(2O^-\) and contains
\(4O^-\), the latter by exponentiation.
There is an exact sequence of \(Q\)-modules
\[
 0\longrightarrow\mathcal T\longrightarrow T
       \xrightarrow{\log}\Lambda\longrightarrow0.
\]
The elements \(\zeta_p^{g^i}\), \(0\le i<2h\), form a
\(\mathbb Z_2\)-basis of \(O\).
Indeed, a polynomial of degree at most \(p-1\) with zero constant
coefficient cannot be a nonzero scalar multiple of \(\Phi_p\).
The differences
\(\zeta_p^{g^i}-\zeta_p^{-g^i}\), \(0\le i<h\), consequently form a
basis of \(O^-\), with \(s\) acting by the signed cyclic shift.
Thus
\[
 O^-\simeq\mathbb Z_2[s]/(s^h+1).
\]
The determinant of \(Q\) on the commensurable lattice \(\Lambda\)
is also \(D\); hence \(Q\) is injective there and
\(|\Lambda/Q\Lambda|=2^{v_2(D)}\).
The exact sequence above gives
\[
 0\longrightarrow\mathcal T/Q\mathcal T
 \longrightarrow T/QT
 \longrightarrow\Lambda/Q\Lambda\longrightarrow0.
\]
This proves \([T:QT]=2^{v_2(D)+t}\); the torsion contribution is
essential.

The kernel of \(s-1:T\to P\) is \(T^s=\{1,-1\}\), because the
\(s\)-fixed subring of \(O\) is \(\mathbb Z_2\).
The operator \(Q\) is an automorphism on this kernel, since \(Q(1)\)
is odd.  Therefore \(s-1\) induces
\[
 T/QT\xrightarrow{\sim}P/QP.
\]
Explicitly, if \((s-1)x=(s-1)Qy\), then
\(x-Qy\in T^s=QT^s\), so \(x\in QT\); surjectivity follows from
the definition of \(P\).

Modulo \(2\), the signed cyclic module becomes
\(\mathbb F_2[X]/(X^h-1)\).  The nullity of multiplication by
\(\overline Q\) on this module is \(b\).
In the Smith form over \(\mathbb Z_2\), exactly \(b\) diagonal entries
of \(Q\) have positive valuation.  Their valuations sum to \(v_2(D)\);
therefore \(v_2(D)\ge b\).

If \(u=1+2a\in T\), its norm-one equation gives
\(a+c(a)+2a\,c(a)=0\).  Thus \(\overline a\in \mathscr M\), and multiplication
of units shows that \(\alpha\) is an equivariant homomorphism.
The elements
\[
 \beta_i=\overline\zeta_p^{\,g^i}
             +\overline\zeta_p^{-g^i},\qquad0\le i<h,
\]
form a basis of \(\mathscr M\), with cyclic \(s\)-action.
For any \(\overline a\in O/2O\), choose a lift \(a\in O\), put
\(v=1+2a\), and set \(u=v/c(v)\).  Then
\(\alpha(u)=\overline a+c(\overline a)\).
This trace map hits every \(\beta_i\), proving that \(\alpha:T\to \mathscr M\)
is onto.
It follows that the induced map on \(Q\)-quotients is onto and
\[
 |\mathscr M/\overline QM|=2^b.
\]
Its source has order \(2^{v_2(D)+t}\), so it is an isomorphism
exactly when \(v_2(D)+t=b\).
The inequality already proved and \(t\ge0\) make this equivalent to
\(t=0\) and \(v_2(D)=b\).

For completeness, the logarithmic convention does not silently lose
the sign contribution.  The logarithm series gives
\[
 \alpha_{\log}(u)=\alpha(u)+\alpha(u)^2=(1+[2])\alpha(u).
\]
Write \(2=g^m\pmod p\).
On \(\mathscr M\), the operator \([2]\) is \(s^m\), and
\(\gcd(h,m)=e\): the order of \(2\) modulo signs is \(d/2\) for even
\(d\), and \(d\) for odd \(d\).
Consequently
\[
 \gcd(X^h-1,X^m-1)=X^e-1.
\]
Multiplication by \(1+X^m\) on \(\mathscr M/\overline QM\) is invertible
exactly when \(t=0\); more precisely its cokernel has dimension \(t\).
When \(t=0\) and \(v_2(D)=b\), both this operator and the map induced
by \(\alpha\) are isomorphisms.  Conversely an isomorphism induced by
\(\alpha_{\log}\) forces equality of source and target orders, and therefore
the same two conditions.
\end{proof}

The class-number identity
\(v_2(D)=v_2(h_p^-)\) makes these conditions an effective test for the
local hypothesis of the global unit criterion.  This corollary
asserts no global arithmetic saturation.

